%/*% pdflatex -file-line-error -interaction nonstopmode %

\documentclass{article}
\usepackage{graphicx}
\graphicspath{ {./images/2.2/} }
\usepackage{amsmath}
\usepackage{amsfonts}
\usepackage{bookmark}
\usepackage{exercises}
\usepackage{amsthm}
\usepackage{yhmath}
\usepackage{placeins}
\newcommand*{\oldneg}{\mathord{\sim}}

\theoremstyle{definition}
\newtheorem{definition}{Definition}

\setlength{\parindent}{.5in}
\setlength{\parskip}{12pt}

\title{Proofs Answers}
\begin{document}
\maketitle

This is a document for proving the theorems that are listed in the textbook.

\section*{2.3}
$(a+b)(a-b)=a^2-b^2$
\[
\begin{array}{rl}
	& (a+b)(a-b) \\
=	& (a+b)a-(a+b)b \\
=	& a^2+ba-ab-b^2 \\
=	& a^2-b^2
\end{array}
\]
\newline
$(a+b)^2=a^2+2ab+b^2$
\[
\begin{array}{rl}
	& (a+b)^2=(a+b)(a+b) \\
=	& (a+b)a+(a+b)b \\
=	& a^2+ab+ab+b^2 \\
=	& a^2+2ab+b^2
\end{array}
\]
\newline
$(a-b)^2=a^2-2ab+b^2$
\[
\begin{array}{rl}
	& (a-b)^2=(a-b)(a-b) \\
=	& (a-b)a-(a-b)b \\
=	& a^2-ab-ab+b^2 \\
=	& a^2-2ab+b^2
\end{array}
\]
\newline
$(a+b)^3=a^3+3a^2b+3ab^2+b^3$
\[
\begin{array}{rl}
	& (a+b)^3=(a+b)(a+b)(a+b) \\
=	& (a^2+2ab+b^2)(a+b) \\
=	& (a^2+2ab+b^2)a+(a^2+2ab+b^2)b \\
=	& a^3+2a^2b+ab^2+a^2b+2ab^2+b^3 \\
=	& a^3+3a^2b+3ab^2+b^3
\end{array}
\]
\newline
$(a-b)^3=a^3-3a^2b+3ab^2-b^3$
\[
\begin{array}{rl}
	& (a-b)^3=(a-b)(a-b)(a-b) \\
=	& (a^2-2ab+b^2)(a-b) \\
=	& (a^2-2ab+b^2)a-(a^2-2ab+b^2)b \\
=	& a^3-2a^2b+ab^2-a^2b+2ab^2-b^3 \\
=	& a^3-3a^2b+3ab^2-b^3
\end{array}
\]

\section*{2.4}
$a^3+b^3=(a+b)(a^2-ab+b^2)$
\[
\begin{array}{rl}
	& (a+b)(a^2-ab+b^2) \\
=	& (a^2-ab+b^2)a+(a^2-ab+b^2)b \\
=	& a^3-a^2b+ab^2+a^2b-ab^2+b^3 \\
=	& a^3+b^3
\end{array}
\]
$a^3-b^3=(a-b)(a^2+ab+b^2)$
\[
\begin{array}{rl}
	& (a-b)(a^2+ab+b^2) \\
=	& (a^2+ab+b^2)a-(a^2+ab+b^2)b \\
=	& a^3+a^2b+b^2a-a^2b-ab^2-b^3 \\
=	& a^3-b^3
\end{array}
\]
$ra+rb+sa+sb=(r+s)(a+b)$
\[ ra+rb+sa+sb=r(a+b)+s(a+b)=(r+s)(a+b) \]

\section*{2.5}

\subsection*{Polynomial Long Division}


We consider a polynomial $f(x)$. It can be equally written as \[ a_nx^n+a_n-1x^{n-1}+...+a_0x^0 \] where $n$ is the degree of $f(x)$. We suppose 
that $f(x)$ is being divided by some $d(x)$. This can be expanded as \[ b_mx^m+b_mx^{m-1}+...+b_0x^0 \] where $m$ is the degree of $d(x)$. Define 
the set of polynomials $p_1(x),p_2(x),...,p_j(x)$ in the following way: $p_1(x)=f(x)-d(x)(\frac{a_n}{b_m})$; 
$p_2(x)=p_1(x)-d(x)(\frac{t}{b_m})$\textrm{ where $t$ is the leading term of $p_1(x)$}; and, in general, $p_i(x)=p_{i-1}(x)-(\frac{t_i}{b_m}$)\textrm{ where 
$t_i$ is the leading term of $p_{i-1}$}. The last term, $p_j(x)$, is such that the degree of $p_j(x)$ is less than $m$. We are to show that 
$\frac{f(x)}{d(x)}=\frac{t_1}{b_m}+\frac{t_2}{b_m}+...+\frac{t_{j-1}}{b_m}+\frac{p_j}{d(x)}$. This can be shown by using double induction, in particular, 
induction on $m$ and $n$. First, we consider $n$, and consider the base case where the degree is $0$. This means that, according to how $p_j(x)$ was 
defined, the only $t_i$ possible would be when $d(x)$ has a degree of $0$ as well, making the result in this case just $\frac{f(x)}{d(x)}$. If 
the degree of $d(x)$ was greater than $0$, this would mean there would be no $p_i$'s, meaning there would only be the remainder, $\frac{f(x)}{d(x)}$, 
completing the proof when $n=0$. Now, suppose that the result holds for some $n$. In the base case where $m=0$, $\frac{t_1}{b_m}=a_{n+1}$. Also, 
every $p_i(x)$ differs from $p_{i-1}$ by only subtracting the leading term of $p_{i-1}$, meaning every $\frac{t_i}{b_m}$ is equal to the corresponding 
$a_i$. There is no remainder in this case (because $m$ is 0), thus completing the base case. Now, suppose that for some $m$, the result holds for $n+1$. 
Now consider the case of $m+1$. Let $d_1(x)$ be $d(x)$ without the term $b_{m+1}$. 

 We can split the sequence \[ \frac{a_{n+1}x^{n+1}+a_nx^n+...+a_0x^0}{d(x)} \] into \[ \frac{a_{n+1}x^{n+1}}{d(x)}+
\frac{a_nx^n+...+a_0x^0}{d(x)} \] by the induction assumption, \[ \frac{a_nx^n+...+a_0x^0}{d(x)}=\frac{t_1}{b_{m+1}}+\frac{t_2}{b_{m+1}}+...+\frac{t_{j-1}}
{b_{m+1}}+\frac{p_j(x)}{d(x)} \]
\end{document}
